% GENERATED FILE. Edit canonical lessons, then run npm run generate:books.
% canonical-source-hash: fnv1a64:f54ac9b9ac3440c7
% canonical-lessons: PA-C110-shakti, PA-C110-hakk, PA-C110-adhikar, PA-C110-mauka, PA-C110-jaidad

\chapter{Power, Right, Authority, Chance, Property}
\label{ch:pa-power-right-authority-chance-property}
\hlchaptermodality{110}

\begin{chapteropening}
\textbf{By the end of this chapter:} \emph{I can say power, a right, authority, a chance and property.}

Complete the last lesson of chapter 110: I can say power, a right, authority, a chance and property.
\end{chapteropening}

\section[\texorpdfstring{\pa{ਸ਼ਕਤੀ} (shaktī)}{ਸ਼ਕਤੀ (shaktī)}]{\pa{ਸ਼ਕਤੀ} (shaktī) — power}

\label{lesson:PA-C110-shakti}



\begin{quote}
Before the new one: say the Punjabi for near, then the Punjabi for strength.
\end{quote}

\subsection*{You'll want to know: \pa{ਸ਼ਕਤੀ}}
\textbf{\pa{ਸ਼ਕਤੀ}} — \emph{shaktī} — \textquotedblleft{}power\textquotedblright{}.

\begin{grammarlens}[title={What you've built}]
One more word for this chapter.
\end{grammarlens}

\subsection*{Your turn}
\begin{itemize}
\raggedright
  \item \emph{Say it:} \emph{shaktī}
  \item \emph{Say it:} \emph{shaktī}, once more
  \item \emph{Say it:} say \emph{shaktī}
  \item \emph{Recall:} say the Punjabi for near, then the Punjabi for strength, then say \emph{shaktī} again
\end{itemize}

\begin{tcolorbox}[breakable,colback=teal!4,colframe=teal!35!black,title={Before you move on}]
What does \pa{ਸ਼ਕਤੀ} mean? (\textquotedblleft{}Power\textquotedblright{}.) Say it once more.
\end{tcolorbox}
\section[\texorpdfstring{\pa{ਹੱਕ} (hakk)}{ਹੱਕ (hakk)}]{\pa{ਹੱਕ} (hakk) — a right}

\label{lesson:PA-C110-hakk}



\begin{quote}
Before the new one: say the Punjabi for strength, then the Punjabi for power.
\end{quote}

\subsection*{You'll want to know: \pa{ਹੱਕ}}
\textbf{\pa{ਹੱਕ}} — \emph{hakk} — \textquotedblleft{}a right\textquotedblright{}.

\begin{grammarlens}[title={What you've built}]
One more word for this chapter.
\end{grammarlens}

\subsection*{Your turn}
\begin{itemize}
\raggedright
  \item \emph{Say it:} \emph{hakk}
  \item \emph{Say it:} \emph{hakk}, once more
  \item \emph{Say it:} say \emph{hakk}
  \item \emph{Recall:} say the Punjabi for strength, then the Punjabi for power, then say \emph{hakk} again
\end{itemize}

\begin{tcolorbox}[breakable,colback=teal!4,colframe=teal!35!black,title={Before you move on}]
What does \pa{ਹੱਕ} mean? (\textquotedblleft{}A right\textquotedblright{}.) Say it once more.
\end{tcolorbox}
\section[\texorpdfstring{\pa{ਅਧਿਕਾਰ} (adhikār)}{ਅਧਿਕਾਰ (adhikār)}]{\pa{ਅਧਿਕਾਰ} (adhikār) — authority, a right}

\label{lesson:PA-C110-adhikar}



\begin{quote}
Before the new one: say the Punjabi for power, then the Punjabi for a right.
\end{quote}

\subsection*{You'll want to know: \pa{ਅਧਿਕਾਰ}}
\textbf{\pa{ਅਧਿਕਾਰ}} — \emph{adhikār} — \textquotedblleft{}authority, a right\textquotedblright{}.

\begin{grammarlens}[title={What you've built}]
One more word for this chapter.
\end{grammarlens}

\subsection*{Your turn}
\begin{itemize}
\raggedright
  \item \emph{Say it:} \emph{adhikār}
  \item \emph{Say it:} \emph{adhikār}, once more
  \item \emph{Say it:} say \emph{adhikār}
  \item \emph{Recall:} say the Punjabi for power, then the Punjabi for a right, then say \emph{adhikār} again
\end{itemize}

\begin{tcolorbox}[breakable,colback=teal!4,colframe=teal!35!black,title={Before you move on}]
What does \pa{ਅਧਿਕਾਰ} mean? (\textquotedblleft{}Authority, a right\textquotedblright{}.) Say it once more.
\end{tcolorbox}
\section[\texorpdfstring{\pa{ਮੌਕਾ} (maukā)}{ਮੌਕਾ (maukā)}]{\pa{ਮੌਕਾ} (maukā) — a chance, an opportunity}

\label{lesson:PA-C110-mauka}



\begin{quote}
Before the new one: say the Punjabi for a right, then the Punjabi for authority.
\end{quote}

\subsection*{You'll want to know: \pa{ਮੌਕਾ}}
\textbf{\pa{ਮੌਕਾ}} — \emph{maukā} — \textquotedblleft{}a chance, an opportunity\textquotedblright{}.

\begin{grammarlens}[title={What you've built}]
One more word for this chapter.
\end{grammarlens}

\subsection*{Your turn}
\begin{itemize}
\raggedright
  \item \emph{Say it:} \emph{maukā}
  \item \emph{Say it:} \emph{maukā}, once more
  \item \emph{Say it:} say \emph{maukā}
  \item \emph{Recall:} say the Punjabi for a right, then the Punjabi for authority, then say \emph{maukā} again
\end{itemize}

\begin{tcolorbox}[breakable,colback=teal!4,colframe=teal!35!black,title={Before you move on}]
What does \pa{ਮੌਕਾ} mean? (\textquotedblleft{}A chance, an opportunity\textquotedblright{}.) Say it once more.
\end{tcolorbox}
\section[\texorpdfstring{\pa{ਜਾਇਦਾਦ} (jāidād)}{ਜਾਇਦਾਦ (jāidād)}]{\pa{ਜਾਇਦਾਦ} (jāidād) — property}

\label{lesson:PA-C110-jaidad}



\begin{quote}
Before the new one: say the Punjabi for authority, then the Punjabi for a chance.
\end{quote}

\subsection*{You'll want to know: \pa{ਜਾਇਦਾਦ}}
\textbf{\pa{ਜਾਇਦਾਦ}} — \emph{jāidād} — \textquotedblleft{}property\textquotedblright{}.

\begin{grammarlens}[title={What you've built}]
One more word for this chapter.
\end{grammarlens}

\subsection*{Your turn}
\begin{itemize}
\raggedright
  \item \emph{Say it:} \emph{jāidād}
  \item \emph{Say it:} \emph{jāidād}, once more
  \item \emph{Say it:} say \emph{jāidād}
  \item \emph{Recall:} say the Punjabi for authority, then the Punjabi for a chance, then say \emph{jāidād} again
\end{itemize}

\begin{tcolorbox}[breakable,colback=teal!4,colframe=teal!35!black,title={Before you move on}]
What does \pa{ਜਾਇਦਾਦ} mean? (\textquotedblleft{}Property\textquotedblright{}.) Say it once more.
\end{tcolorbox}
